% LuaLaTeX document — compile with: lualatex moment_t9_no_eshelby.tex  (twice)
\documentclass[11pt,a4paper]{article}

\usepackage{fontspec}
\setmainfont{Latin Modern Roman}
\setmonofont{Latin Modern Mono}

\usepackage{amsmath,amssymb,amsthm,bm}
\usepackage{booktabs}
\usepackage{geometry}
\geometry{margin=2.5cm}
\usepackage{hyperref}
\hypersetup{colorlinks=true,linkcolor=blue!60!black,citecolor=green!50!black,urlcolor=blue!70!black}
\usepackage{enumitem}
\usepackage{xcolor}
\usepackage{mathtools}

\newcommand{\D}{\Delta}
\newcommand{\order}{\mathcal{O}}
\newcommand{\dd}{\mathrm{d}}
\newcommand{\Gz}{G^{0}}
\newcommand{\half}{\tfrac{1}{2}}

\title{The nine-component single-site T-matrix\\ from moment integrals,
       without Eshelby}
\author{}
\date{}

\begin{document}
\maketitle

\begin{abstract}
The cubic single-site T-matrix is usually obtained by importing Eshelby's
uniform-eigenstrain result and patching it for a cube. This note derives it
instead from the moment integrals of the elastodynamic Green's tensor over the
cube. Eshelby is never invoked: it \emph{emerges} as the first-gradient
truncation of a closed hierarchy, and the derivation says exactly what that
truncation costs. The route is constructive throughout --- every integral is in
closed form, the $r=0$ distributional content is carried automatically rather
than patched in, and the resulting $9\times 9$ block reduces under $O_h$ to
three scalars. Two independent implementations agree on those scalars to
$4\times10^{-16}$.
\end{abstract}

\section{The closed set}

Let the scatterer be a cube $V=[-\D/2,\D/2]^3$ of contrast
$\delta c_{ijkl}=\delta\lambda\,\delta_{ij}\delta_{kl}
+\delta\mu(\delta_{ik}\delta_{jl}+\delta_{il}\delta_{jk})$ and $\delta\rho$,
embedded in an isotropic background with Green's tensor $\Gz_{ij}$. Expanding
the internal field about the cube centre and projecting the
Lippmann--Schwinger equation onto the monomials $1$, $x_p$, $x_px_q$ gives a
closed system in the three unknowns $u_i(0)$, $\partial_pu_i(0)$,
$\partial_p\partial_qu_i(0)$:
\begin{align}
\bigl(\delta_{ij}-\omega^2\delta\rho\,G_{ij}\bigr)u_j(0)
  &+\bigl(N^{r}_{in,k}\,\delta c_{nksj}
  -\half\omega^2\delta\rho\,K^{rs}_{ij}\bigr)\partial_r\partial_su_j(0)
  = u^0_i(0), \label{eq:A11}\\[2pt]
\bigl(\delta_{pr}\delta_{ij}+\omega^2\delta\rho\,N^{r}_{ij,p}
  &-M_{in,pk}\,\delta c_{nkrj}\bigr)\partial_ru_j(0)
  = \partial_pu^0_i(0), \label{eq:A22}\\[2pt]
-\omega^2\delta\rho\,M_{ij,pq}u_j(0)
  &+\bigl(\delta_{qr}\delta_{ps}\delta_{ij}+Q^{r}_{in,pqk}\,\delta c_{nksj}
  -\half\omega^2\delta\rho\,P^{rs}_{ij,pq}\bigr)\partial_r\partial_su_j(0)
  = \partial_p\partial_qu^0_i(0). \label{eq:A33}
\end{align}
The six coefficient tensors are moments of $\Gz$ over the cube,
\begin{equation}
\begin{aligned}
G_{ij}&=\int_V \Gz_{ij}\,\dd V, &
K^{rs}_{ij}&=\int_V \Gz_{ij}\,x_rx_s\,\dd V, &
N^{r}_{in,k}&=\int_V \partial_k\Gz_{in}\,x_r\,\dd V,\\
M_{in,pk}&=\int_V \partial_p\partial_k\Gz_{in}\,\dd V, &
P^{rs}_{ij,pq}&=\int_V \partial_q\partial_p\Gz_{ij}\,x_rx_s\,\dd V, &
Q^{r}_{in,pqk}&=\int_V \partial_q\partial_p\partial_k\Gz_{in}\,x_r\,\dd V,
\end{aligned}
\label{eq:moments}
\end{equation}
graded by $(D,W)$, the number of derivatives and the degree of the $x$-weight:
\[
G\,(0,0)\quad K\,(0,2)\quad N\,(1,1)\quad M\,(2,0)\quad P\,(2,2)\quad Q\,(3,1).
\]

\paragraph{A convention that is easy to get backwards.} The derivatives in
\eqref{eq:moments} are with respect to the \emph{source} point. Since $\Gz$
depends on the separation, $\partial' = -\partial$ and a moment with $D$
derivatives differs from its field-derivative value by $(-1)^D$. Only $N$ and
$Q$ are affected. This is invisible to every symmetry, parity and scaling
check, so it must be fixed by convention rather than discovered by testing.

\section{The parity decoupling}

$D+W$ is even for five of the six moments and odd only for $N$. Equivalently,
$u$, $\partial u$ and $\partial\partial u$ carry degrees $0$, $1$, $2$, and
even couples only to even. The consequence is
\eqref{eq:A22}: \emph{it involves nothing else}. The first gradient is
determined by its own block, with no reference to $u$ or to
$\partial\partial u$.

This is the structural fact the rest of the note rests on. The leading
scattering from a modulus contrast is the stress dipole
$\int_V \delta c:\nabla u\,\dd V$, which needs only $\partial u$. Therefore
\emph{no amount of work on the second-gradient block
\eqref{eq:A33} can change the leading far field}. Whatever \eqref{eq:A33}
contributes reaches the far field only through the $\order((ka)^2)$ coupling
in \eqref{eq:A11}.

\section{The moments, in closed form}

Write the elastodynamic tensor as
\begin{equation}
\Gz_{in}(\bm r)=\frac{1}{4\pi\mu}\Bigl[\delta_{in}\,g_\beta(r)
+\frac{1}{k_\beta^2}\,\partial_i\partial_n\bigl(g_\beta(r)-g_\alpha(r)\bigr)\Bigr],
\qquad g_k(r)=\frac{e^{ikr}}{r}.
\label{eq:Gdecomp}
\end{equation}
Its static limit reproduces Kupradze exactly, with
$a_0=(\lambda+3\mu)/8\pi\mu(\lambda+2\mu)$ and
$b_0=(\lambda+\mu)/8\pi\mu(\lambda+2\mu)$. Expanding $g_k$ in powers of $r$
turns every moment in \eqref{eq:moments} into a finite sum of \emph{scalar}
moments
\begin{equation}
E\bigl[m;d_1\ldots d_D;w_1\ldots w_W\bigr]
 =\Bigl\langle \partial_{d_1}\!\cdots\partial_{d_D} r^m,\;
   x_{w_1}\!\cdots x_{w_W}\,\mathbf{1}_V\Bigr\rangle .
\label{eq:E}
\end{equation}

\paragraph{Why these are distributions, not integrals.} For $m=-1$ the
classical integrand is not absolutely integrable: $\partial^4(1/r)$ is a
harmonic, hence traceless, tensor of degree $-5$, so its angular average
vanishes and the radial divergence is multiplied by zero. The value then
depends on the shape of the exclusion. This is precisely the delta-function
term that naive quadrature misses --- and gets the \emph{sign} wrong without.

The remedy is not to patch a delta in afterwards but to pair against the test
function $w\,\mathbf{1}_V$ from the start. Using
$\partial_q\mathbf{1}_V=-n_q\delta_{\partial V}$,
\begin{equation}
\bigl\langle \partial_q T',\,w\mathbf{1}_V\bigr\rangle
 =\oint_{\partial V} n_q\,w\,T'\,\dd A
 -\bigl\langle T',\,(\partial_qw)\mathbf{1}_V\bigr\rangle ,
\label{eq:peel}
\end{equation}
with $T'=\partial_{\text{rest}}r^m$. Two features make \eqref{eq:peel} exact
rather than merely plausible: $T'$ is evaluated \emph{only on the boundary},
at distance $\D/2$ from the origin, where every derivative of $r^m$ is smooth
and bounded; and the recursion lowers both $D$ and $W$, terminating on an
ordinary convergent integral. The delta content is inside the surface term and
cannot be dropped. The sharp check is
\begin{equation}
\sum_p E[-1;\{p,p\};\,] = -4\pi ,
\label{eq:lap}
\end{equation}
which a spherical excision returns as $0$.

\paragraph{Conditioning.} The integrals emit logarithms of Pell units,
$(2\pm\sqrt3)^n=A_n\pm B_n\sqrt3$ solving $x^2-3y^2=1$. Left unreduced these
are numerically lethal: $\log(72010600134783751-41575339372323900\sqrt3)$ is a
$35$-digit cancellation between integers of order $10^{17}$ and evaluates to
\texttt{Indeterminate} in double precision. Reduced by $\log(A_n-B_n\sqrt3)=n\log(2-\sqrt3)$
it is exact to machine precision. Every closed form below is Pell-reduced.

\section{The nine-component block}

Specialise \eqref{eq:A22}. With $\delta\rho=0$ --- the regime in which the
single-site modulus T-matrix is defined --- the block is
\begin{equation}
A^{(2)}_{(p,i),(r,j)}
 =\delta_{pr}\delta_{ij}-M_{in,pk}\,\delta c_{nkrj},
\label{eq:A22static}
\end{equation}
a $9\times 9$ operator built from $M$ alone. It contains neither $Q$ nor $P$:
the second-gradient moments do not appear.

Under $O_h$ the nine components decompose as
$\mathbf{3}\otimes\mathbf{3}=A_{1g}\oplus E_g\oplus T_{1g}\oplus T_{2g}$, and
\eqref{eq:A22static} is diagonal in that basis. The strain concentration in
each channel is the inverse of the corresponding eigenvalue:
\begin{align}
A_{1g}\ \text{(bulk)}: &\quad
  1+\frac{3\delta\lambda+2\delta\mu}{3(\lambda+2\mu)}, \label{eq:bulk}\\[4pt]
T_{2g}\ \text{(off-diagonal shear)}: &\quad
  1+2\,\delta\mu\,S_{\text{shear}},\qquad
  S_{\text{shear}}=\frac{\pi(\lambda+2\mu)-\sqrt3(\lambda+\mu)}
                        {3\pi\mu(\lambda+2\mu)}, \label{eq:t2g}\\[4pt]
E_{g}\ \text{(diagonal shear)}: &\quad
  1+2\,\delta\mu\,S_{\text{diag}},\qquad
  S_{\text{diag}}=\frac{3\sqrt3(\lambda+\mu)+2\mu\pi}
                        {6\pi\mu(\lambda+2\mu)} . \label{eq:eg}
\end{align}

That $E_g\neq T_{2g}$ is not a discrepancy: they are inequivalent
representations, and their difference \emph{is} the cubic anisotropy, carried
elsewhere in this project as $\D\mu^*_{\text{diag}}-\D\mu^*_{\text{off}}$. A
sphere would have no such split; a cube must.

\section{Where Eshelby went}

Eshelby's uniformity theorem says the internal field of an inclusion under
uniform remote strain is uniform \emph{if and only if} the inclusion is an
ellipsoid. The usual construction imports the ellipsoidal result and adapts
it. Nothing of the sort happens here. Equations
\eqref{eq:bulk}--\eqref{eq:eg} follow from the moment integrals and the group
theory alone.

What \emph{does} correspond to Eshelby is the truncation: retaining only
$\partial u$ and discarding $\partial\partial u$ is exactly the
uniform-internal-strain ansatz. So the classical result is recovered as the
first-gradient closure of \eqref{eq:A11}--\eqref{eq:A33}, and the hierarchy
makes the cost of that closure explicit rather than hidden.

\paragraph{And the cost does not shrink by going further.} Because
\eqref{eq:A22} is self-contained, adding the second-gradient block cannot
improve the leading far field. This is not an estimate; it is a consequence of
the parity decoupling. It also explains an otherwise puzzling observation: the
tier sequence $T_9\to T_{27}\to T_{57}$ in the shear channel runs
$1.000\to0.7525\to0.9692$ --- oscillating, not converging. A polynomial
expansion about the centre cannot represent the edge and corner stress
singularities of a cube at any order, so the series is asymptotic. Improving
the far field requires changing the \emph{basis}, not adding terms to it.

\section{Validation}

The construction is checked against objects it does not contain.

\begin{center}
\begin{tabular}{@{}lll@{}}
\toprule
Check & Against & Result \\
\midrule
$\sum_p E[-1;\{p,p\}]=-4\pi$ & analytic & exact \\
$\nabla^2 r=2/r$, $\nabla^2r^3=12r$, $\nabla^2r^5=30r^3$ & analytic & exact \\
$\int_V \dd V/r$ & solid-angle quadrature & $3.4\times10^{-9}$ \\
$M$ components $A,B,C$ & independent first-principles route & $10$ digits \\
$M$, all $81$ components & 2012 archive closed form & exact \\
$T_{2}^{c}$, $T_{3}^{c}$ & production \texttt{effective\_contrasts} & $3.6\times10^{-16}$,
$1.3\times10^{-15}$ \\
\bottomrule
\end{tabular}
\end{center}

The last row is the load-bearing one. The production code reaches the same
quantities by a completely different route --- tabulated master integrals with
an explicit Eshelby delta correction --- and the two agree to machine
precision. Matching the identification
\[
\text{amp}^{\text{off}}_e=\frac{1}{1-2T_2},\qquad
\text{amp}^{\text{diag}}_e=\frac{1}{1-2T_2-T_3}
\]
against \eqref{eq:t2g}--\eqref{eq:eg} predicts, with no free parameters,
\[
T_2^{c}=-\delta\mu\,S_{\text{shear}},\qquad
T_3^{c}=2\,\delta\mu\,\bigl(S_{\text{shear}}-S_{\text{diag}}\bigr),
\]
and both hold to $10^{-15}$. Two independent implementations agreeing is the
evidence standard; one implementation agreeing with itself is not.

\section{What this does and does not establish}

\textbf{Established.} The nine-component single-site T-matrix follows from the
moment integrals with no ellipsoidal input; the moments are exact and carry
their distributional content correctly; the block reduces to three closed-form
scalars; and those scalars are confirmed by an independent implementation.

\textbf{Not established.} That the result is \emph{accurate} for a cube. It is
the first-gradient truncation of an asymptotic hierarchy, and
\eqref{eq:bulk}--\eqref{eq:eg} inherit whatever error that truncation carries.
The construction makes the approximation explicit and bounds where it can
live --- in the internal-field basis, not in the integrals --- but it does not
remove it. Improving on it means replacing the polynomial ansatz, for instance
by supplying the three channel concentrations from an independent exact
computation while keeping everything else in this note intact.

\end{document}
